\begin{figure}[!htbp]
\centering
\begin{tikzpicture}[
  box/.style={font=\small,align=center,text width=4.2cm,minimum height=1.04cm,
              draw=#1,rounded corners=2pt,inner sep=3pt,fill=#1!4},
  labeltext/.style={font=\footnotesize,align=center,inner sep=3pt},
  arrow/.style={-{Stealth[length=1.8mm]},draw=ink,line width=.55pt}]
\node[box=sea] (in) at (0,0)
  {Inputs in $\mathcal T_S$\\$Y_\nu^{m_\nu}=\lambda_\nu$};
\node[box=sea] (out) at (8.4,0)
  {Output in $\mathcal T_S$\\$Y_\nu^{m_\nu}=\lambda_\nu$};
\node[box=ink] (lift) at (0,-2.23)
  {Lifted coefficients in $E[t]$\\$\deg P_a,\deg Q_a<r$};
\node[box=ink] (recover) at (8.4,-2.23)
  {Product coefficients in $E[t]$\\$\deg R_a\le2r-2+|S|$};
\node[box=rust] (splitin) at (0,-4.46)
  {Inputs over $E[t]/(t^H)$\\$Y_\nu^{m_\nu}=1$};
\node[box=rust] (splitout) at (8.4,-4.46)
  {Product over $E[t]/(t^H)$\\$Y_\nu^{m_\nu}=1$};
\draw[arrow] (in) -- node[labeltext,right] {undo the unit scaling\\lift with $z_1=t-1$} (lift);
\draw[arrow] (lift) -- node[labeltext,right] {reduce modulo $t^H$\\apply the unit scaling} (splitin);
\draw[arrow] (splitin) -- node[labeltext,above] {Fourier evaluation}
                        node[labeltext,below] {$W_S$ jet products\\and interpolation} (splitout);
\draw[arrow] (splitout) -- node[labeltext,left] {undo the unit scaling\\read $R_a(t)$ in $E[t]$} (recover);
\draw[arrow] (recover) -- node[labeltext,left] {substitute $t=1+z_1$\\truncate and scale} (out);
\end{tikzpicture}
\caption{Multiplication in a binomial block using one parameter for all
axes. The upper row uses length-$r$ jets at $t=1$ and the lower row
uses length-$H$ jets at $t=0$. Each axis wraps at most once in a
monomial product, so $H=2r+d$ captures every coefficient of the
ordinary polynomial in the middle row. This degree bound justifies
reading the polynomial before evaluating it at the other endpoint.}
\label{poly:fig:deformation}
\end{figure}
